\appendix

\section{The exact computation}\label{app:exact}

The program \texttt{chl\_coefficients.py} of the archive \cite{Bha26c} proves Theorem~\ref{thm:bounded} in exact integer arithmetic, in five steps.
\begin{enumerate}
\item[(B1)] \emph{Orbits.} A charge $(m,n,\ell)$ corresponds to the binary quadratic form $f(s,t)=ms^2+2\ell st+2nt^2$ of discriminant $-8D$, and $\Gamma_1(2)$ acts by $f\mapsto f\circ g$. The orbits with given $D$ are enumerated through the Gauss-reduced forms of discriminant $-8D$, each with a class of bottom rows modulo $2$, taken modulo the automorphs of the reduced form.
\item[(B2)] \emph{A representative in $\mathcal A$.} In every orbit the program finds a charge with $0\le\ell<2m$ satisfying \eqref{eq:setA}, which it tests in the form $2(jm+\ell m_1)^2+m_1^2\cdot2D\ge2m^2$ for $m_1\in2\ZZ\setminus\{0\}$ and odd $j$ with $m_1\mid(j^2-1)/4$, the range of $(m_1,j)$ being finite because the left side grows.
\item[(B3)] \emph{Coefficients.} The coefficients $c_{\mathcal R}(m,n,\ell)$ are computed from \eqref{eq:product} by the recursion for the exponential of \eqref{eq:Ll}, with polynomial products in FLINT \cite{Har10}. By Lemma~\ref{lem:chamber}(ii) each charge is read in the cheaper of its two orientations.
\item[(B4)] \emph{Signs.} For every orbit, $d^*=(-1)^{\ell+1}c_{\mathcal R}(m,n,\ell)$ is computed for every representative found in the box, the values are compared, and the sign is checked. At $m=1,2$ the values are compared with Propositions~\ref{prop:Q1closed} and~\ref{prop:Q2closed}.
\item[(B5)] \emph{Least representatives.} On the whole range, every charge of least $Q^2$ in its orbit is checked to lie in $\mathcal A$, as Lemma~\ref{lem:leastA} predicts.
\end{enumerate}
For $D\le\Dexact$ there are $\Norbits$ orbits, with $Q^2\le\Qmaxexact$ for the least representative; all have $d^*>0$. The smallest value relative to the extreme term occurs at small $D$, where $x$ is close to $1$; this is the region in which the leading-family error of Lemma~\ref{lem:mainerr} is largest and which Theorem~\ref{thm:bounded} covers for $Q^2\le7$.

\section{Interval computations}\label{app:interval}

All the numerical statements used in the proofs are certified in ball arithmetic with Arb \cite{Joh17} through python-flint; each program prints its enclosures.
\begin{enumerate}
\item[(I1)] \texttt{pmin.py}: Proposition~\ref{prop:pmin}. Product truncated at $n=60$ with an explicit tail enclosure, $4001$ points on the half circle, Lipschitz constant $L=2P(r_0)r_0P'(r_0)$.
\item[(I2)] \texttt{gmono2.py}: Lemma~\ref{lem:gfacts}. Monotonicity of $\phi$ on $[3,60]$ by mean-value enclosures on an adaptive cover; the range $[60,\infty)$ through the closed form \eqref{eq:halfint}.
\item[(I3)] \texttt{cert.py}: Theorem~\ref{thm:cert}. For each $m$ the program lists the leading family $\mathcal L_m$, the walls with $rs\ge2$ and the parity walls through their parameters $(r,s,a,b)$, encloses $t(r,s,a,b)$, and evaluates the seven bounds of Section~\ref{sec:errors}. In \eqref{eq:XT} it uses $32$ in place of $16$.
\item[(I4)] \texttt{analytic.py}, \texttt{analytic2.py}: Theorem~\ref{thm:analytic}.
\item[(I5)] \texttt{polar.py}: the wall sums \eqref{eq:P0}, \eqref{eq:Pinf}, used also for the checks (c), (d) of Section~\ref{ssec:evidence}.
\end{enumerate}

\section{The certified margins}\label{app:margins}

Table~\ref{tab:margins} lists, for $3\le Q^2\le100$, the total bound $\epsilon(m,x_0(m))$ of Theorem~\ref{thm:cert} and the margin $0.53816-\epsilon$.

\begin{table}[ht]
\centering
\footnotesize
\setlength{\tabcolsep}{3.5pt}
\caption{The certified total error $\epsilon(m,x_0(m))$ of Theorem~\ref{thm:cert} and the margin $0.53816-\epsilon$, for $m=Q^2$.}\label{tab:margins}
\begin{tabular}{rll|rll|rll|rll}
\hline
$m$ & $\epsilon$ & margin & $m$ & $\epsilon$ & margin & $m$ & $\epsilon$ & margin & $m$ & $\epsilon$ & margin \\
\hline
3 & $6.04\cdot10^{-6}$ & 0.5380 & 28 & 0.1540 & 0.3842 & 53 & 0.0913 & 0.4467 & 78 & 0.0648 & 0.4732 \\
4 & $5.47\cdot10^{-5}$ & 0.5379 & 29 & 0.1500 & 0.3883 & 54 & 0.0898 & 0.4482 & 79 & 0.0640 & 0.4740 \\
5 & $2.87\cdot10^{-4}$ & 0.5377 & 30 & 0.1460 & 0.3921 & 55 & 0.0884 & 0.4496 & 80 & 0.0633 & 0.4747 \\
6 & $9.79\cdot10^{-4}$ & 0.5370 & 31 & 0.1420 & 0.3958 & 56 & 0.0870 & 0.4510 & 81 & 0.0626 & 0.4754 \\
7 & 0.0024 & 0.5356 & 32 & 0.1390 & 0.3993 & 57 & 0.0857 & 0.4523 & 82 & 0.0619 & 0.4761 \\
8 & 0.3460 & 0.1917 & 33 & 0.1350 & 0.4026 & 58 & 0.0844 & 0.4536 & 83 & 0.0612 & 0.4768 \\
9 & 0.3350 & 0.2030 & 34 & 0.1320 & 0.4058 & 59 & 0.0831 & 0.4549 & 84 & 0.0605 & 0.4775 \\
10 & 0.2930 & 0.2450 & 35 & 0.1290 & 0.4089 & 60 & 0.0819 & 0.4561 & 85 & 0.0599 & 0.4781 \\
11 & 0.2790 & 0.2594 & 36 & 0.1260 & 0.4118 & 61 & 0.0807 & 0.4573 & 86 & 0.0592 & 0.4788 \\
12 & 0.2670 & 0.2714 & 37 & 0.1230 & 0.4145 & 62 & 0.0795 & 0.4585 & 87 & 0.0586 & 0.4794 \\
13 & 0.2550 & 0.2828 & 38 & 0.1210 & 0.4172 & 63 & 0.0784 & 0.4596 & 88 & 0.0580 & 0.4800 \\
14 & 0.2450 & 0.2930 & 39 & 0.1180 & 0.4197 & 64 & 0.0773 & 0.4607 & 89 & 0.0574 & 0.4806 \\
15 & 0.2350 & 0.3026 & 40 & 0.1160 & 0.4222 & 65 & 0.0763 & 0.4617 & 90 & 0.0568 & 0.4812 \\
16 & 0.2260 & 0.3115 & 41 & 0.1130 & 0.4245 & 66 & 0.0753 & 0.4627 & 91 & 0.0563 & 0.4817 \\
17 & 0.2180 & 0.3199 & 42 & 0.1110 & 0.4268 & 67 & 0.0742 & 0.4638 & 92 & 0.0557 & 0.4823 \\
18 & 0.2100 & 0.3277 & 43 & 0.1090 & 0.4289 & 68 & 0.0733 & 0.4647 & 93 & 0.0551 & 0.4829 \\
19 & 0.2030 & 0.3350 & 44 & 0.1070 & 0.4310 & 69 & 0.0723 & 0.4657 & 94 & 0.0546 & 0.4834 \\
20 & 0.1960 & 0.3419 & 45 & 0.1050 & 0.4330 & 70 & 0.0714 & 0.4666 & 95 & 0.0541 & 0.4839 \\
21 & 0.1900 & 0.3483 & 46 & 0.1030 & 0.4350 & 71 & 0.0705 & 0.4675 & 96 & 0.0536 & 0.4844 \\
22 & 0.1840 & 0.3544 & 47 & 0.1010 & 0.4368 & 72 & 0.0696 & 0.4684 & 97 & 0.0530 & 0.4850 \\
23 & 0.1780 & 0.3601 & 48 & 0.0994 & 0.4386 & 73 & 0.0688 & 0.4692 & 98 & 0.0525 & 0.4855 \\
24 & 0.1730 & 0.3655 & 49 & 0.0976 & 0.4404 & 74 & 0.0679 & 0.4701 & 99 & 0.0520 & 0.4860 \\
25 & 0.1670 & 0.3706 & 50 & 0.0960 & 0.4420 & 75 & 0.0671 & 0.4709 & 100 & 0.0516 & 0.4864 \\
26 & 0.1630 & 0.3754 & 51 & 0.0943 & 0.4437 & 76 & 0.0663 & 0.4717 &  & &  \\
27 & 0.1580 & 0.3799 & 52 & 0.0928 & 0.4452 & 77 & 0.0655 & 0.4725 &  & &  \\
\hline
\end{tabular}
\end{table}

